\documentclass[11pt]{book}
\usepackage[utf8]{inputenc}
\usepackage[T1]{fontenc}
\usepackage{geometry}
\geometry{a5paper, margin=1.5cm}
\usepackage{parskip}
\usepackage{enumitem}
\usepackage{booktabs}
\usepackage{array}
\usepackage{titlesec}
\usepackage{fancyhdr}
\usepackage{charter}

\pagestyle{fancy}
\fancyhf{}
\fancyhead[LE,RO]{\thepage}
\fancyhead[RE]{\textit{The Artificer's Codex}}
\fancyhead[LO]{\textit{Korrin Iron-Stroke}}
\renewcommand{\headrulewidth}{0.4pt}

\titleformat{\chapter}[hang]{\normalfont\Huge\bfseries}{}{0pt}{}
\titleformat{\section}{\normalfont\Large\bfseries}{}{0pt}{}
\titleformat{\subsection}{\normalfont\itshape}{}{0pt}{}

\setcounter{secnumdepth}{3}

\begin{document}

\thispagestyle{empty}
\vspace*{2cm}
\begin{center}
    {\Huge \textbf{The Artificer's Codex}}
    \vspace{0.5cm}
    {\Large A Guide to Permanent Enchantment and Iterative Craft}
    \vspace{1cm}

    \textit{``The Monad teaches that nothing is ever finished — only briefly sufficient. A blade that cannot be improved is a blade that has already begun to die.''}
    \vspace{0.5cm}

    --- \textbf{Korrin Iron-Stroke}, \\
    Runekeeper of the Clockwork Monad, \\
    Master Artificer of Khaz-Vurim
    \vfill

    \texttt{Khaz-Vurim, 832 AR}
\end{center}

\newpage
\section*{Preface: A Letter from Korrin}

\emph{To whoever finds this battered journal ---}

I am Korrin Iron-Stroke. I am an Aeler, which means I count my breaths, I measure my words, and I never, ever pour the ninth cup. I am a Runekeeper of the Clockwork Monad, which means I have spent my life learning that perfection is not a destination --- it is a direction. And I am a craftsman, which means I understand that a well-made thing is a promise, and a promise is a debt, and a debt is a weight that must be measured before it is incurred.

This book is not a manual. It is a ledger. It records what I have learned about making things that last --- not by accident, but by design. I have made blades that remember the hands that held them. I have forged bells that ring only when a promise is true. I have crafted keys that open doors that were never meant to open, and I have sealed doors that should never be opened again.

I have also made mistakes. I have created things that I should not have. I have bound obligations I could not pay. I have seen what happens when a craftsman forgets that every tool is a creditor, and every creation is a covenant.

The Clockwork Monad taught me this: \textbf{A thing that is made is a thing that is owed.} The moment you shape a piece of metal, you have incurred a debt to the shape you are creating. The moment you pour a spell into steel, you have bound yourself to that spell's nature. The moment you craft an item of power, you have entered into a contract --- with the item, with the Patron who whispered the design, and with the world that will remember what you have made.

So before you begin, ask yourself: \textbf{What am I willing to owe?}

\vspace{0.5cm}
\hfill --- Korrin Iron-Stroke \\
\hfill Khaz-Vurim, 832 AR \\
\hfill (stamped with the seal of the Iteration Chorus)

\newpage
\tableofcontents

\newpage
\chapter{The Philosophy of Making}

\section{A Thing That is Made is a Thing That is Owed}

The world runs on obligation. Everything you touch, everything you shape, everything you improve --- it all remembers the hand that made it. In the Aeler holds, we say: ``The mountain remembers every keystone.'' In the Monad's forge, we say: ``The machine remembers every iteration.''

When you craft a permanent magic item, you are not simply creating a tool. You are creating a creditor. The item will demand something from you --- time, attention, sacrifice, or Obligation. The Patron who guided your work will expect a return on their investment. And the world itself will note the new shape you have forced into existence.

\textbf{This is not a burden to be feared.} It is a relationship to be managed. A well-crafted item is like a well-sworn oath: it serves you faithfully if you serve it faithfully. But a neglected item --- like a broken promise --- will turn against you.

\section{The Spectrum of Power}

Not all magic items are created equal. I have learned to measure the weight of a creation before I begin the work. Use this scale as a guide:

\begin{table}[h]
\centering
\begin{tabular}{@{}llllll@{}}
\toprule
\textbf{Tier} & \textbf{XP Cost} & \textbf{DV} & \textbf{Timer} & \textbf{Obligation} & \textbf{Example} \\
\midrule
Cantrip-Infused & 2 XP & 2 & 4 segments & 1 segment & A knife that always returns, a bell that rings at dawn \\
Minor Enchantment & 4 XP & 3 & 4 segments & 2 segments & A cloak that hides your face, a ring that warns of lies \\
Standard Item & 8 XP & 4 & 6 segments & 4 segments & A sword that burns with frost, a shield that remembers oaths \\
Major Relic & 12 XP & 5 & 8 segments & 6 segments & A crown that bends wills, a staff that speaks to stones \\
Legendary Artifact & 16+ XP & 6+ & 10+ segments & 8+ segments & The Kalrantha Amulet, the Ninth Ledger, the Hammer That Unmakes \\
\bottomrule
\end{tabular}
\end{table}

\vspace{0.5cm}
\textbf{Design Note:} The Obligation cost is not optional. Even if you have no Patron, the item itself creates a debt --- to the world, to the materials, to the memory of what you have shaped. A character who crafts a Major Relic without a Patron may find that the Relic itself \emph{becomes} a Patron, demanding service and attention.

\newpage
\chapter{The Tools of the Trade}

\section{The Workshop --- A Place of Covenants}

Before you can craft anything of worth, you must have a workshop. A workshop is not a room; it is a threshold --- a place where the ordinary becomes the extraordinary, and where the world's rules are bent by the craftsman's will.

\begin{table}[h]
\centering
\begin{tabular}{@{}llll@{}}
\toprule
\textbf{Tier} & \textbf{XP Cost} & \textbf{Downtime to Establish} & \textbf{Benefits} \\
\midrule
Basic Forge & 4 XP & 1 day & Craft Minor items; +1 die to Craft rolls \\
Enchanted Workshop & 8 XP & 1 week & Craft Standard items; +2 dice to Craft rolls; reduce crafting DV by 1 \\
Monadic Forge & 12 XP & 1 month & Craft Major items; +3 dice to Craft rolls; reduce crafting DV by 2; can craft items that generate Obligation to the Clockwork Monad \\
\bottomrule
\end{tabular}
\end{table}

\vspace{0.5cm}
\textbf{The Ninth Rule of the Workshop:} Every workshop must have a ninth shelf, a ninth peg, or a ninth drawer that is never used. The space is for the Hollow. The Hollow is patient. The Hollow watches.

\section{Components --- The Language of Materials}

A magic item is not made from steel and leather alone. It is made from moments. The materials matter because they carry memory. A sword forged from iron that was once a plowshare remembers the soil. A crystal cut from a stone that once held a river remembers the water. A hide tanned from a beast that once killed remembers the hunt.

\begin{table}[h]
\centering
\begin{tabular}{@{}lll@{}}
\toprule
\textbf{Component Type} & \textbf{Examples} & \textbf{DV Adjustment} \\
\midrule
Common & Steel, leather, basic herbs, copper, simple dyes & --- \\
Uncommon & Uncommon herbs, monster teeth, blessed salt, moonlit iron & +1 DV \\
Rare & Dragon scale, phoenix feather, storm-caught lightning, Ykrul war-braid & +2 DV \\
Legendary & True name of a dead god, heart of a living mountain, a tear from a Veiled Sister & +3 DV \\
\bottomrule
\end{tabular}
\end{table}

\textbf{The Ninth Component:} A component that has been touched by the Hollow --- a stone from a bridge that has no ninth arch, a bell that has never rung, a coin that was never spent. Using such a component reduces the DV by 1, but adds 1 Obligation to the Hollow.

\section{The Role of Obligation}

When you craft a permanent magic item, you incur Obligation. This is not a punishment --- it is the price of power.

\textbf{Obligation from Crafting:}
\begin{itemize}
    \item \textbf{Cantrip-Infused:} 1 segment to the Clockwork Monad (or the item itself)
    \item \textbf{Minor Enchantment:} 2 segments
    \item \textbf{Standard Item:} 4 segments
    \item \textbf{Major Relic:} 6 segments
    \item \textbf{Legendary Artifact:} 8+ segments
\end{itemize}

\textbf{Clearing Crafting Obligation:}
\begin{itemize}
    \item \textbf{Maintenance:} Spend 1 XP or a downtime scene maintaining the item.
    \item \textbf{Iteration:} Improve the item (spend XP to upgrade it). Each upgrade clears 1 Obligation.
    \item \textbf{Service:} Perform a service for the Monad or the item's spirit.
    \item \textbf{Delegation:} A Cap 4+ follower can manage upkeep for you (see Chapter 4).
\end{itemize}

\newpage
\chapter{The Crafting Process}

\section{The Five Steps of Creation}

\begin{enumerate}
    \item \textbf{Conceive} --- Design the item. What does it do? What does it demand? What is its cost?
    \item \textbf{Gather} --- Acquire the components. This is often a quest in itself.
    \item \textbf{Shape} --- Work the materials. This is the crafting roll, using a timer.
    \item \textbf{Imbue} --- Pour the power into the item. This is the final test.
    \item \textbf{Bind} --- Seal the work. Pay the XP, incur the Obligation, and witness the item's birth.
\end{enumerate}

\section{The Timer --- Patience is a Metal}

Each crafting project has a timer. The timer represents the iterative process --- the failures, the refinements, the moments where you step back and see the flaw you missed.

\begin{table}[h]
\centering
\begin{tabular}{@{}ll@{}}
\toprule
\textbf{Item Tier} & \textbf{Timer Size} \\
\midrule
Cantrip-Infused & 4 segments \\
Minor Enchantment & 4 segments \\
Standard Item & 6 segments \\
Major Relic & 8 segments \\
Legendary Artifact & 10+ segments \\
\bottomrule
\end{tabular}
\end{table}

\textbf{Ticking the Timer:}
\begin{itemize}
    \item \textbf{Success on a Craft roll:} Advance 1 segment.
    \item \textbf{Partial Success:} Advance 1 segment, but the GM gains 1 SB (complications).
    \item \textbf{Miss:} No advance. The GM may tick a Complication Timer.
\end{itemize}

\textbf{When the Timer Fills:} The item is ready for the Imbuement roll. If you have rare components, you may add them now to reduce the DV of the Imbuement roll.

\section{The Imbuement Roll}

The Imbuement roll is the moment of truth. Roll Wits + Arcana (or appropriate skill) against the item's DV.

\begin{itemize}
    \item \textbf{Clean Success:} The item is complete. Perfect. No flaws.
    \item \textbf{Success with SB:} The item is complete, but the GM gains 2 SB to spend on a later complication --- a rival who senses the power, a flaw that emerges in a year, a debt that the item will demand.
    \item \textbf{Partial:} The item is complete, but with a flaw (see Table 3.1). You may choose to accept the flaw or start over.
    \item \textbf{Miss:} The item is not complete. The components are consumed. You may try again, but you must acquire new components.
\end{itemize}

\section{Flaws --- The Price of Imperfection}

A flawed item is not a failure. It is a story.

\begin{table}[h]
\centering
\begin{tabular}{@{}lll@{}}
\toprule
\textbf{d6} & \textbf{Flaw} & \textbf{Mechanical Effect} \\
\midrule
1 & Fragile & The item breaks on a critical failure. Can be repaired with a downtime scene. \\
2 & Cursed & The item has a minor curse (e.g., attracts insects, makes you smell of brimstone, whispers secrets). \\
3 & Unstable & The item's effect is unpredictable; roll 1d6, on a 1 it backfires. \\
4 & Hungry & The item demands sacrifice --- blood, time, or a memory --- once per session. \\
5 & Bound to Patron & The item is tied to a specific Patron; using it against the Patron's interests triggers backlash. \\
6 & The Ninth Flaw & The item has a secret flaw that only the Hollow knows. The GM may reveal it at a dramatically appropriate moment. \\
\bottomrule
\end{tabular}
\end{table}

\newpage
\chapter{Maintaining Your Creations}

\section{Upkeep --- The Cost of Keeping}

Every magic item requires upkeep. Neglect a blade, and it will dull. Neglect a ring, and it will forget its power. Neglect a pact, and the creditor will collect.

\begin{table}[h]
\centering
\begin{tabular}{@{}lll@{}}
\toprule
\textbf{Upkeep Method} & \textbf{Cost} & \textbf{Effect} \\
\midrule
Efficient & XP equal to the item's cost divided by 3 (minimum 1) & Minimal effort. The item remains Functional. \\
Intensive & 1 XP + a downtime scene & Significant personal involvement. The item may gain a minor benefit (e.g., +1 die on its next use). \\
Delegated & A Cap 4+ follower manages the item & The follower pays the upkeep (no XP cost). Risk: if the follower's Loyalty drops, all delegated items become Neglected. \\
\bottomrule
\end{tabular}
\end{table}

\section{Item Status}

\begin{table}[h]
\centering
\begin{tabular}{@{}ll@{}}
\toprule
\textbf{Status} & \textbf{Effect} \\
\midrule
Functional & The item works as intended. \\
Neglected & --1 die to the item's effects. Can be restored with a downtime scene or 1 XP. \\
Compromised & The item's effects are unreliable; roll 1d6, on a 1 the effect fails or backfires. Requires a quest or 4 XP to restore. \\
\bottomrule
\end{tabular}
\end{table}

\newpage
\chapter{Sample Creations}

\section{Cantrip-Infused Items (2 XP, DV 2, Timer 4)}

\begin{itemize}
    \item \textbf{Returning Knife:} A knife that returns to your hand when thrown. It does not fly; it simply appears in your palm after a heartbeat.
    \item \textbf{Dawn Bell:} A bell that rings once at dawn, wherever it is. It cannot be silenced.
    \item \textbf{Truth Coin:} A coin that lands on a specific face when a lie is told within Near range.
\end{itemize}

\section{Minor Enchantments (4 XP, DV 3, Timer 4)}

\begin{itemize}
    \item \textbf{Emberfang Dagger:} Glows faintly when undead are near; once per scene, gain +1 die to detect hidden spirits.
    \item \textbf{Cloak of the Unseen:} +1 die to Stealth in dim light; once per scene, reroll a failed Stealth check.
    \item \textbf{Vial of the Last Breath:} Heal 1 Fatigue when drunk. One use.
\end{itemize}

\section{Standard Items (8 XP, DV 4, Timer 6)}

\begin{itemize}
    \item \textbf{Storm Lantern:} Once per session, as an action, create a gust of wind that extinguishes all mundane flames in Near and gives +1 Position to disengage.
    \item \textbf{Pendant of the Witness:} Once per session, when you detect a lie, the liar suffers 1 Fatigue (no save).
    \item \textbf{Raven's Key:} Once per day, unlock any mundane lock or open any door that is not magically sealed.
\end{itemize}

\section{Major Relics (12 XP, DV 5, Timer 8)}

\begin{itemize}
    \item \textbf{Crown of Unspoken Oaths:} Once per day, compel a target to answer one question truthfully (Resist DV 4). The target does not remember answering.
    \item \textbf{Staff of Walking Stones:} Command up to three stones to move as if they were alive. They can create cover, block passages, or carry messages.
\end{itemize}

\section{Legendary Artifacts (16+ XP, DV 6+, Timer 10+)}

\begin{itemize}
    \item \textbf{The Kalrantha Amulet (see Chapter 6)}
    \item \textbf{The Ninth Ledger (see Chapter 6)}
    \item \textbf{The Hammer That Unmakes:} Once per arc, strike a surface and shatter any non-living object (size up to a building). The Hammer demands a memory in return.
\end{itemize}

\newpage
\chapter{Artifacts --- Items That Owe You}

\section{What Makes an Artifact?}

An artifact is not crafted --- it is discovered, inherited, or earned. It has a will of its own, however faint. It remembers every hand that held it, every oath spoken over it, every debt left unpaid.

An artifact has three special properties:

\begin{enumerate}
    \item \textbf{Obligation.} When you claim an artifact, you mark Obligation to the artifact itself. This Obligation is tracked like Patron Obligation (Capacity = Spirit + Presence).
    \item \textbf{Rites.} An artifact grants one or two specific Rites that you can invoke without a Patron, a Codex, or a Symbol. Using a Rite from an artifact costs the normal Obligation.
    \item \textbf{Indestructible (mostly).} If an artifact is lost or destroyed, it reappears at the next dawn --- somewhere nearby, often in the hand of the one who owes it the most Obligation.
\end{enumerate}

\section{Sample Artifacts}

\begin{table}[h]
\centering
\begin{tabular}{@{}lll@{}}
\toprule
\textbf{Name} & \textbf{Obligation} & \textbf{Effect} \\
\midrule
The Copper Bell of Unspoken Truths & 2 & Once per session, ring the bell. The next lie told within Near range is heard as a discordant chime; you gain +1 die to expose it. \\
Ledger of the Debtor Saint & 3 & You may record a single promise in the ledger. While the promise is unfulfilled, you gain +1 die to all rolls that directly advance it. When the promise is fulfilled, clear 1 Obligation. \\
The Mourner's Cloak & 1 & When you pull the hood up, you become invisible to anyone who is actively grieving (GM's judgment). The invisibility lasts until you speak or attack. \\
Hearthstone of the Lost Hearth & 2 & Once per day, set the stone on a threshold. For the next hour, that threshold cannot be crossed by anyone with hostile intent (Resolve DV 4 to resist). \\
\bottomrule
\end{tabular}
\end{table}

\newpage
\chapter{The Iterative Art --- Upgrading Items}

\section{The Principle of Iteration}

The Clockwork Monad teaches that nothing is ever finished. A sword can always be sharper. A shield can always be lighter. A ring can always remember more.

\textbf{Upgrading an Item:}
\begin{itemize}
    \item Spend XP equal to the difference between the old tier and the new tier.
    \item Roll Craft + Wits (DV = new tier + 1).
    \item \textbf{Success:} The item improves. Clear 1 Obligation.
    \item \textbf{Partial:} The item improves, but gains a flaw.
    \item \textbf{Miss:} The item does not improve. You may try again.
\end{itemize}

\section{Example Upgrades}

\begin{table}[h]
\centering
\begin{tabular}{@{}llll@{}}
\toprule
\textbf{Starting Item} & \textbf{Upgrade} & \textbf{Cost} & \textbf{New Effect} \\
\midrule
Emberfang Dagger (Minor) & $\to$ Standard & 4 XP & Glows when undead are near; once per scene, gain +2 dice to detect hidden spirits. \\
Cloak of the Unseen (Minor) & $\to$ Standard & 4 XP & +2 dice to Stealth in dim light; once per scene, reroll a failed Stealth check; once per session, become invisible for one exchange. \\
Storm Lantern (Standard) & $\to$ Major & 4 XP & Once per session, create a gust of wind; once per day, call a small storm that lasts one scene. \\
\bottomrule
\end{tabular}
\end{table}

\newpage
\chapter{The Ninth Rule}

The Ninth Rule is simple: \textbf{Never make the ninth item.}

I have seen it happen. A craftsman makes eight swords. Each is better than the last. The ninth would be perfect. The Monad whispers: ``Just one more iteration. Just one more pass. Just one more refinement.''

The ninth item is always perfect. It is also always hungry. It remembers what you sacrificed to make it --- the materials, the time, the sleepless nights, the friendships you neglected, the promises you broke to finish it. It knows your weaknesses. And it will use them.

I have seen a master smith forge a blade that could cut through any armor. The ninth blade. He did not survive the year. The blade found his heart when he was sleeping.

I have seen an archivist create a tome that could answer any question. The ninth book. She forgot her own name. The book remembered. The book is still waiting.

So when you have made eight items, stop. Count your breaths. Pour eight cups. Walk eight steps. Do not make the ninth.

The Monad is patient. It can wait for the tenth.

\newpage
\appendix
\chapter{Sample Components by Region}

\begin{table}[h]
\centering
\begin{tabular}{@{}lllll@{}}
\toprule
\textbf{Region} & \textbf{Common} & \textbf{Uncommon} & \textbf{Rare} & \textbf{Legendary} \\
\midrule
Aeler & Iron, copper, breath-glass & Deep-crystal, bell-bronze & Breath-dust, under-vein ore & Heart of the mountain, the Ninth Seal \\
Acasia & Rusted steel, old wood & Hedge-witch salt, broken milestones & Curse-iron, Pale Causeway stone & A bridge's memory, a broken emperor's crown \\
Vhasia & Tournament steel, wine-stained cloth & A sun-cracked crown, a knight's gorget & Fenwood ash, Lence marble & The shattered sun's shard, a queen's promise \\
Thepyrgos & Ink, parchment, candle wax & Synod seal wax, a bell's clapper & A stair from the Unfinished Stair, a Lamp-Lighter's oil & A page from the Ninth Ledger, a fossil stair \\
Ykrul & Horse-hair, leather, bone & Wolf-tooth, steppe-salt, hostage string & A Khagan's braid, a white squall's eye & The Hoof-Print, a horse's true name \\
Silkstrand & Silk, ink, bridge-stone & Dye from the canals, a weaver's thread & A Matron's seal, a bridge's memory & The ninth bridge's stone, the Ragman's favor \\
\bottomrule
\end{tabular}
\end{table}

\newpage
\chapter{The Artificer's Mantras}

\begin{itemize}
    \item ``Measure twice. Forge once. Iterate always.''
    \item ``The eighth is mercy. The ninth is the Hollow's.''
    \item ``A tool is not a servant. It is a witness.''
    \item ``Perfection is not a destination. It is a debt.''
    \item ``The Monad does not forgive. It reschedules.''
    \item ``When you craft, you incur. When you incur, you owe. When you owe, you are bound.''
    \item ``The Hollow is the interest on every promise you have ever made.''
    \item ``Eight is enough. Eight is peace. Eight is the number of the living.''
    \item ``The ninth item is always perfect. It is also always hungry.''
\end{itemize}

\newpage
\section*{Colophon}

\emph{This book was written in the workshop of Korrin Iron-Stroke, Khaz-Vurim, in the Year of the Iron Harvest. The ink is lamp-black and sweat. The binding is iron and leather. The ninth page is blank --- it is for the Hollow, should it ever wish to read.}

\vspace{1cm}
\hfill --- Korrin Iron-Stroke \\
\hfill Runekeeper of the Clockwork Monad \\
\hfill Master Artificer of the Aeler Holds

\vspace{0.5cm}
\hfill \emph{``The coin that never spends is the one you don't remember taking.''}

\end{document}
